\section*{AI 协助}
\label{sec:ai-assistance}

我们使用了基于 LLM 的协助（包括 OpenAI Codex）来起草、重构和编辑手稿；准备中文阅读版；检索相关文献与书目信息；以及协助分析与绘图代码。生成式图像工具用于创建和修改解释性框架图。

AI 协助还参与了架构的信息论解释、数学主张及其论证的表述，以及实验方法与结果的讨论。这包括渐进式实体对齐的处理、检索——证据分离，以及依赖查询的读取预算。这些论证以论文中声明的假设为条件；AI 协助不构成独立的证明验证或实证校验。

本论文准备流程中未创建任何合成评测数据集。报告的基准分数取自有记录的实验运行，而不是生成的性能估计。作者对最终文本、引用、数学主张、代码、图表和实证结论负责，包括对 AI 协助材料的核验。
